\documentclass{article}
\usepackage[utf8]{inputenc}
\usepackage[T1]{fontenc}
\usepackage[margin=2cm]{geometry}
\usepackage{hyperref}
\sloppy
\begin{document}
\section*{Historical source}
Encyclopaedia Britannica scan, source row 2569. Parse 3.5 extracted the text. Super 3.5 VL corrected it from an enlarged crop of each block. It can still contain recognition errors.
Page headers, footers, and page numbers are removed.

\subsection*{ABYSSINIA (E0249)}
\textit{Continues an entry from an earlier page.}

Abyssinia. her affairs entirely to Frumentius, who made use of his influence to diffuse the Christian religion throughout the country, and at last was appointed bishop of Axuma. It is said, however, that the court and principal people, if not the nation in general, relapsed into idolatry, which continued to prevail till the year 521, when they were again converted by their king Adad or Aidog.

The two things are these princes it remarkable for an expedition into Arabia Felix, called by the Mohammedan writers the war of the elephant, and which was undertaken on the following occasion: The temple of Mecca, situated nearly in the middle of the Arabian peninsula, had been held in the greatest veneration for near 1400 years; probably from the notion entertained by the people in the neighbourhood, that Adam pitched his tent on that spot. Here also was a black stone supposed to possess extraordinary sanctity, as being that on which Jacob laid his head when he had the vision of angels. The most probable account of the real origin of this temple, according to Mr Bruce, is, that it was built by Sesostris, and that he himself was worshipped there under the name of Osiris.

On account of the veneration in which this tower and idol were held by the Arabians, Mr Bruce supposes that the thought was first suggested of making it the emporium of the trade between India and Africa; but Abra, in order to divert it into another channel, built a very large temple near the Indian ocean in the country of the Homerites; and, to encourage the resort of people to this new temple, he bestowed upon it all the privileges of the former which stood in the city of Mecca. The tribe of Arabians named Koreish, in whose country Mecca stood, being exceedingly alarmed at the thoughts of having their temple deserted, entered the new one in the night, burned all that could be consumed, and besmeared the remains with human excrements. Abra, provoked at this sacrilege, assembled a considerable army, with which he

The termination of the war, according to the Arabian historians, was miraculous. A vast number of birds named Ababil came from the sea, having faces like lions; each carrying in its claws a small stone about the size of a pea, which they let fall upon the Ethiopian army in such numbers, that every one of them was destroyed. At this time it is said that the small-pox first made its appearance; and the more probable account of the destruction of the Ethiopian army is, that they perished by this distemper.

The war of the elephant is supposed to have terminated in the manner above mentioned about the year 360; from which time to that of Elesbaan, named also Caleb, and probably the same with the Adad or A-

dag already mentioned, we meet with nothing remarkable in the Ethiopic history. He engaged in a war with the Homerites or Sabeans in Arabia Felix, whom he overthrew in battle, and put an end to their kingdom; after which he embraced the Christian religion in token of gratitude for the success he had met with. In the time of this prince a violent persecution of the Christians took place in Arabia. The Jewish religion had now spread itself far into that peninsula; and in many places the professors of it were become absolute masters of the country, insomuch that several Jewish principalities had been erected, the sovereigns of which commenced a severe persecution against the Christians. Among the rest, one Phineas distinguished himself by his cruelty, having prepared a great number of furnaces or pits filled with fire, into which he threw those who refused to renounce Christianity.

The Christians applied for relief to the emperor Justin; but he being at that time engaged in a war with the Persians, could not interfere; however, in the year 522, he sent an embassy to Elesbaan, who was now also a member of the Greek church, intreating him to exert himself for the relief of the Christians of Arabia. On this the emperor commanded his general Abreha, governor of the Arabian province Yemen, to march to the assistance of Aretas, son to a prince of the same name whom Phineas had burnt; while he himself prepared to follow with a more considerable force. But before the arrival of the Ethiopian monarch, young Aretas had marched against Phineas, and entirely defeated him. In a short time afterwards the emperor himself arrived, and gave Phineas a second defeat; but notwithstanding these misfortunes, it does not appear that either the principality of Phineas or any of the other Jewish ones, was at this time overturned, though it seems to be certain, that at the time we speak of, the Ethiopians possessed part of the Arabian peninsula. According to the Arabian historians, the war of the elephant, with the miraculous destruction of the Ethiopian army, already mentioned, took place in the reign of Elesbaan.

Some historians mention, that the Ethiopian monarchs embraced the doctrines of Mahomet soon after the impostor made his appearance; but this seems not to be well-founded; though it is certain that the Najashi or Ethiopian governor of Yemen embraced Mahometanism, and that he was related to the royal family. On this occasion, however, the Ethiopians lost all the footing they once had in Arabia; the governors being expelled by Mahomet and his successors. They fled to the African side of the Red sea with numbers of their subjects, where they erected several small king-driven out doms, as Adel, Wypo, Hadea, Mara, and others, of Arabia, which still continue.

During the conquests of the caliphs, the Jews were for some time everywhere driven out of their dominions, or oppressed to such a degree that they voluntarily left them. Ethiopia offered them an asylum: Number of Jews in Ethiopia increased. and in this country they became so powerful, that a revolution in savour of Judaism seemed ready to take place. One family had always preserved an indepen-

dent sovereignty on a mountain called Samen, the royal residence being on the top of a high rock; and several other high and rugged mountains were used by that people as natural fortresses. Becoming by de- grees


\section*{Source record}
\begin{description}
\item[Scan ID] \texttt{britannica:309f765b1ecc83e0a7f7762e44b4a927a7555443:train:2569}
\item[Image file] \texttt{page-2569-e8e9334880f9.jpg}
\item[Image SHA-256] \texttt{e8e9334880f9e63008622bc2d1196536177dfc1b0a45b3abf304dbcf2a4e72e7}
\item[Dataset] NationalLibraryOfScotland/encyclopaedia\_britannica\_illustrated at revision \texttt{309f765b1ecc83e0a7f7762e44b4a927a7555443} (CC0-1.0)
\item[Extraction] model \texttt{nvidia/nemotron-3.5-parse}, raw reply \texttt{extract/raw/page-2569-e8e9334880f9.json}
\item[Raw reply SHA-256] \texttt{4a0681380c4f0844028cf4985bc116fb435e53ca3fd0d5fb46ed8ce07a91e03a}
\item[Correction] model \texttt{nvidia/Nemotron-3.5-Super-VL}, block crops, raw replies \texttt{correct/raw/crop\_nothink/page-2569-e8e9334880f9-b*.json}; super\_vl: 10
\item[Running header] ABY S SINIA. 57
\end{description}
\end{document}
